\subsection{Hypocontinuous bilinear mappings}

In what follows, we shall define a notion which is intermediate between that of a continuous bilinear mapping and that of a separately continuous bilinear mapping.

PROPOSITION 3. --- Let E, F, G be three locally convex spaces, S a family of bounded subsets of E. Let u be a separately continuous bilinear mapping from E × F into G. The following properties are equivalent :

\begin{enumerate}
\def\labelenumi{\alph{enumi})}
\item
  For every neighbourhood \(\mathbf{W}\) of 0 in \(\mathbf{G}\) and every set \(\mathbf{M} \in \mathfrak{S}\) , there exists a neighbourhood \(\mathbf{V}\) of 0 in \(\mathbf{F}\) such that \(u(\mathbf{M} \times \mathbf{V}) \subset \mathbf{W}\) .
\item
  For every set \(\mathbf{M} \in \mathfrak{S}\) , the image of \(\mathbf{M}\) under the mapping \(x \mapsto u(x, .)\) is an equicontinuous subset of \(\mathcal{L}(\mathrm{F}; \mathrm{G})\) .
\item
  The mapping \(y \mapsto u(., y)\) from \(\mathbf{F}\) into \(\mathcal{L}_{\mathfrak{S}}(\mathbf{E}; \mathbf{G})\) is continuous.
\item
  expresses that \(y \mapsto u(., y)\) is continuous at the point 0, on account of the definition of neighbourhoods of 0 in \(\mathcal{L}_{\mathfrak{S}}(\mathrm{E}; \mathrm{G})\) (III, p.~13); likewise a) expresses that the image of M under the mapping \(x \mapsto u(x, .)\) is equicontinuous at the point 0 (III, p.~16).
\end{enumerate}

DEFINITION 2. --- Let u be a bilinear mapping from E × F into G. We say that u is S-hypocontinuous if u is separately continuous and if it verifies one of the equivalent conditions a), b), c) of prop. 3.

The condition \(c)\) of prop. 3 shows that the notion of \(\mathfrak{S}\) -hypocontinuous bilinear mapping depends on \(\mathfrak{S}\) only through the \(\mathfrak{S}\) -topology on \(\mathcal{L}(\mathrm{E},\mathrm{G})\) .

For every set T of bounded subsets of F, we define similarly the notion of T-hypo-continuous mapping, by interchanging the roles of E and F in prop. 3. A separately continuous bilinear mapping u is said to be \((\mathfrak{S}, \mathfrak{T})\) -hypocontinuous if it is both S-hypocontinuous and T-hypocontinuous.

Every continuous bilinear mapping from \(E \times F\) into G is \((\mathfrak{S}, \mathfrak{T})\) -hypocontinuous for every pair \((\mathfrak{S}, \mathfrak{T})\) of sets of bounded subsets: for every neighbourhood W of 0 in G, there exists a neighbourhood U of 0 in E and a neighbourhood V of 0 in F such that \(u(U \times V) \subset W\) ; since every set \(M \in S\) is bounded, there exists \(\lambda > 0\) such that \(\lambda M \subset V\) , and so

\[
u (\mathbf {M} \times \lambda \mathbf {V}) = u (\lambda \mathbf {M} \times \mathbf {V}) \subset u (\mathbf {U} \times \mathbf {V}) \subset \mathbf {W}.
\]

The converse is in general false (III, p.~47, exerc. 3).

PROPOSITION 4. --- Let u be a S-hypocontinuous bilinear mapping from E × F into G. For every set M ∈ S, the restriction of u to M × F is continuous, and u(M × Q) is bounded in G for every bounded subset Q of F.

The first assertion follows from cor. 3 of GT, X, § 2, No.~1. Let W be a neighbourhood of 0 in G; there exists, by hypothesis, a neighbourhood V of 0 in F such that \(u(\mathbf{M} \times \mathbf{V}) \subset \mathbf{W}\) . Since there exists \(\lambda \neq 0\) such that \(\lambda \mathbf{Q} \subset \mathbf{V}\) , we have \(\lambda u(\mathbf{M} \times \mathbf{Q}) = u(\mathbf{M} \times \lambda \mathbf{Q}) \subset \mathbf{W}\) , and this proves the second part of the proposition.

PROPOSITION 5. --- Let u be a (S, T)-hypocontinuous bilinear mapping from E × F into G. For every pair of sets M ∈ S, N ∈ T, u is uniformly continuous on M × N. The proposition follows immediately from prop. 2 of GT, X, § 2, No.~1 and prop. 5 of GT, X, § 2, No.~2.

PROPOSITION 6. --- If F is a barrelled space, every separately continuous bilinear mapping u from E × F into a locally convex space G is S-hypocontinuous for every set S of bounded subsets of E.

In other words, the linear mapping \(y \mapsto u(., y)\) from \(F\) into \(\mathcal{L}_b(E; G)\) is continuous. It is enough (III, p.~30, prop. 3) to prove that the image of every bounded subset \(M\) of \(E\) under \(x \mapsto u(x, .)\) is equicontinuous in \(\mathcal{L}(F; G)\) . But, by virtue of prop. 1 (III, p.~28) this image is a simply bounded subset of \(\mathcal{L}(F; G)\) , and since \(F\) is barrelled, every simply bounded subset of \(\mathcal{L}(F; G)\) is equicontinuous (III, p.~25, th. 1).

Remark. --- Suppose the topology of F is the finest locally convex topology on F for which the linear mappings \(h_{\alpha}:F_{\alpha}\to F\) are continuous (II, p.~27). Then condition c) of prop. 3 (III, p.~30) shows that if E and G are locally convex, then the bilinear mapping \(u:E\times F\to G\) is S-hypocontinuous if and only if each of the bilinear mappings

\[
(x, y _ {\alpha}) \mapsto u (x, h _ {\alpha} (y _ {\alpha}))
\]

from \(\mathbf{E} \times \mathbf{F}_{\alpha}\) into \(\mathbf{G}\) is \(\mathfrak{S}\) -hypocontinuous.

Now suppose that \(\mathbf{E}\) is a locally convex space which is the strict inductive limit of an increasing sequence \((\mathbf{E}_n)\) of closed vector subspaces of \(\mathbf{E}\) (II, p.~33); then every set \(\mathbf{M} \in \mathfrak{S}\) is contained in one of the \(\mathbf{E}_n\) and is bounded in this subspace (III, p.~5, prop. 6). We denote by \(\mathfrak{S}_n\) the family of all subsets belonging to \(\mathfrak{S}\) contained in \(\mathbf{E}_n\) .

Condition a) of prop. 3 (III, p.~30) shows that for a bilinear mapping \(u: \mathrm{E} \times \mathrm{F} \to \mathrm{G}\) to be \(\mathfrak{S}\) -hypocontinuous, it is necessary and sufficient that each of the restrictions \(u_n: \mathrm{E}_n \times \mathrm{F} \to \mathrm{G}\) of \(u\) is \(\mathfrak{S}_n\) -hypocontinuous.

